\begin{tabular}{>{\raggedright\arraybackslash}p{4.2cm}rlc>{\raggedright\arraybackslash}p{5.2cm}}
\toprule
Input & Base case & Range & PSA & Source
\\
\midrule
Surgical circumcision at permanent health facilities, per procedure & \$36.30 & \$24.18--\$45.37 & Beta-PERT & Larson et al. 2015 (MUWRP fixed sites, central Uganda). Low: surgical cost in Duffy et al. 2013. High: +25\% \\
Surgical circumcision at outreach camps, per procedure & \$65.13 & \$48.84--\$76.87 & Beta-PERT & Larson et al. 2015 (MUWRP mobile clinic). High: more remote areas in the same study. Low: -25\% \\
Device circumcision at permanent health facilities, per procedure & \$44.73 & \$32.62--\$49.00 & Beta-PERT & Fixed-site surgery (Larson et al. 2015) plus the PrePex premium of \$7.90 in Duffy et al. 2013. Low: Duffy's \$30.55. High: 35\% above surgery \\
Cost of delivering an incentive payment & \$0.00 & \$0.00--\$2.00 & Fixed & Assumption: administering cash at the facility \\
Circumcisions per HIV infection averted & 12 & 8--20 & Beta-PERT & Kripke et al. 2016, Uganda DMPPT 2.0, ages 20-34 (12-13), 2014-2028. Range: assumption \\
Antiretroviral therapy (ART), cost per person per year & \$522.95 & \$368.27--\$683.94 & Beta-PERT & Kripke et al. 2016 (Uganda HIV investment case, 2013). Range: assumption \\
Years from circumcision to start of ART (averted infection) & 5 & 2--10 & Beta-PERT & Assumption \\
Years on ART per averted infection & 20 & 10--30 & Beta-PERT & Assumption \\
Discount rate for averted treatment costs & 3\% & 1\%--5\% & Fixed & Kripke et al. 2016 used 3\% \\
\bottomrule
\end{tabular}
